% sec-09-additional-information.tex - Additional Information (full IND)
\section{Additional Information}

This section addresses the categories of additional information enumerated for an
original Investigational New Drug (IND) submission under 21 CFR \S 312.23(a)(10),
namely drug dependence and abuse potential, radioactive drugs, pediatric studies,
and any other information bearing on the safe and ethical conduct of the proposed
clinical investigation. Because this is a combined IND and Investigational Device
Exemption (IDE) first-in-human study of orally administered daraxonrasib
(RMC-6236) delivered alongside an on-premises large language model (LLM) directed
eight-arm robotic pancreaticoduodenectomy, the Other Information subsection
additionally consolidates the Physical AI governance, cybersecurity, record
retention, and three-tier oversight content that the autonomy-bearing
investigational device requires for Phase 1 acceptance. Each statement below is
grounded in the protocol (v1.0.0, DOI 10.5281/zenodo.20780121) and is consistent
with FDA Phase 1 guidance \cite{phase1guidance} and with the regulatory adaptations
developed for Physical AI investigations \cite{cfr312paiuni,cfr050paiuni,ichpaistduni}.

\subsection{Drug Dependence and Abuse Potential}

Daraxonrasib has no known dependence or abuse potential and is not a controlled
substance under the Controlled Substances Act. The investigational drug is an oral,
once-daily RAS(ON) multi-selective (pan-RAS) small-molecule inhibitor that binds
the guanosine triphosphate (GTP) bound active form of mutant RAS in a complex with
cyclophilin A and thereby suppresses downstream effector signaling
\cite{daraxwhipple,onpremwhippl}. Its pharmacological mechanism is targeted
oncogenic-signaling inhibition; it has no activity at the central nervous system
receptors, transporters, or ion channels that mediate reinforcement, euphoria,
sedation, or physical dependence, and it does not belong to a pharmacological class
associated with recreational use, tolerance with escalating self-administration, or
a withdrawal syndrome. The drug carries an FDA Breakthrough Therapy designation
(June 2025) for previously treated metastatic pancreatic ductal adenocarcinoma
(PDAC) with KRAS G12 mutations, and the prior and ongoing human experience
summarized in the Previous Human Experience section has not identified any signal
of abuse, dependence, drug-seeking behavior, or diversion.

Accordingly, no abuse-liability program is proposed for this IND. No dedicated
nonclinical abuse-potential studies (for example, self-administration, drug
discrimination, or physical-dependence and withdrawal studies under the FDA
abuse-potential assessment framework) are warranted, and no controlled-substance
scheduling recommendation applies. The investigational product is dispensed under
the accountability controls described in the Chemistry, Manufacturing, and Control
Information and Proposed Clinical Research sections, with drug accountability logs
maintained at the single academic site, so that the limited quantities used for the
planned cohorts (up to eighteen treated participants across three dose levels) are
reconciled from receipt through administration, return, and destruction. Should any
unanticipated central-nervous-system finding suggestive of abuse potential emerge
during the conduct of the study, it would be captured through the routine adverse
event and serious adverse event reporting pathways under 21 CFR \S 312.32 and
reviewed by the Data Safety Monitoring Board (DSMB).

\subsection{Radioactive Drugs}

This subsection is not applicable. Daraxonrasib is not a radioactive drug and
contains no radionuclide. The investigational product is a non-radioactive oral
small-molecule tablet administered once daily; it is neither a positron-emitting
nor a single-photon-emitting radiopharmaceutical, and it is not the subject of a
Radioactive Drug Research Committee (RDRC) determination under 21 CFR \S 361.1. No
radiation dosimetry, no radiation safety committee review, and no radioactive
drug-specific manufacturing, handling, storage, or waste-disposal controls are
required for the investigational agent.

Imaging performed during the study (for example, diagnostic computed tomography and
any standard-of-care nuclear medicine staging that may precede enrollment) is part
of routine clinical care and is not an investigational radioactive drug administered
under this IND; such imaging is performed and interpreted according to institutional
standards and is outside the scope of this subsection. The intraoperative imaging,
hemostasis, and drain placement of operative phase P8 are optical and instrumental
and involve no administered radioactivity. No information under 21 CFR
\S 312.23(a)(10)(ii) is therefore required.

\subsection{Pediatric Studies}

The study population is restricted to adults aged eighteen years or older. Eligible
participants have histologically or cytologically confirmed PDAC that is resectable
or borderline-resectable, a documented KRAS G12 mutation (G12D, G12V, or G12R), an
Eastern Cooperative Oncology Group (ECOG) performance status of 0 to 1, and are
candidates for pancreaticoduodenectomy (Whipple resection), consistent with the
broad pan-KRAS eligibility framework of the RASolute 302 program
\cite{daraxwhipple}. No participant under the age of eighteen will be enrolled,
consented, screened, dosed, or operated upon under this protocol, and pediatric
assent provisions are therefore inapplicable in practice even though the
informed-consent architecture (21 CFR part 50 with ICH E6(R3)) provides for assent
where it would otherwise apply \cite{cfr050paiuni}.

A first-in-human Phase 1 evaluation in a pediatric population is neither warranted
nor proposed at this stage, and a pediatric study plan is deferred. The deferral
rationale rests on three considerations. First, the disease under study, PDAC, is
overwhelmingly a disease of older adults; pancreatic ductal adenocarcinoma in
children is exceedingly rare, the five-year overall survival of PDAC is below
thirteen percent, and the condition is the third leading cause of cancer death in
the United States and the fourth in the European Union \cite{pdac060s2030}, so the
target indication does not arise materially in the pediatric population. Second,
this is a first-in-human study of an autonomy-bearing investigational device
combined with an investigational dosing regimen; the prudent and ethical sequence
is to establish the device and procedure serious adverse event profile, the maximum
tolerated dose and recommended Phase 2 dose, and the technical feasibility of the
LLM-directed robotic Whipple in adults before any consideration of younger
populations, and full autonomy is in any case prohibited at this stage under 21 CFR
\S 312.21(e) \cite{cfr312paiuni}. Third, the staged-autonomy model of this trial
class confines the present study to clinician-controlled teleoperation and
supervised operation (Stage 1), reserving higher-autonomy operation for future
study; extension to a pediatric population, if ever pursued, would follow
successful completion of the adult program and would be the subject of a separate
pediatric study plan with its own age-appropriate formulation, dosing, device
sizing, and assent and parental-permission provisions.

Because this is an investigation of a serious and life-threatening oncology
indication that does not occur meaningfully in children, the study qualifies for
the conditions under which a pediatric assessment is not required for an adult-only
first-in-human Phase 1 program; any subsequent pediatric commitments would be
addressed at the appropriate later stage of development in consultation with the
FDA. No pediatric data are submitted with this original IND.

\subsection{Other Information}

The autonomy-bearing investigational device requires governance, cybersecurity,
record-retention, and oversight controls beyond those of a conventional small-molecule
Phase 1 study. This subsection consolidates that Physical AI overlay. The
on-premises LLM operates strictly as a second-opinion advisory oracle: it is
hash-pinned to a deterministic seed and a repository commit, isolated outside the
vendor kinematic stack, and never granted control authority over the robotic arms;
full autonomy is prohibited under 21 CFR \S 312.21(e) \cite{cfr312paiuni}, and every
proposed robot-patient command passes a verification-before-generation ten-gate
suite returning accept, block, or escalate, aligned with the Verification Before
Generation in Physical AI Oncology Trials Act of 2026 (H.R. 9510)
\cite{congress9510}. The governance design follows the trust and accountability
principles developed for clinical Physical AI \cite{clintrustpai,paiautospons}.

\subsubsection{Physical AI Governance and the Three-Tier Oversight Structure}

Independent safety oversight operates on three coordinated tiers layered over the
reviewing Institutional Review Board (IRB) and the Sponsor-Investigator. The DSMB,
composed of members independent of the conduct of the study with surgical,
oncologic, biostatistical, and device-safety expertise, reviews accumulating safety
data at each cohort boundary and at any triggered review and holds binding halt
authority. The Independent Safety Monitor (ISM) is designated for each procedure and
provides real-time, case-level safety oversight independent of the operating team,
with the authority to call an immediate stop. The Physical AI Safety Review
Committee meets on a cadence of ninety days or less (and ad hoc on any triggering
event) to review the autonomy-bearing behavior of the system, including the
telemetry and audit record, the verification-before-generation gate-surface
decisions, the Unified Safety Level (USL) readiness rating, and any software change,
and recommends USL reassessment or oversight changes to the Sponsor-Investigator.
The three tiers are coordinated so that no autonomy-bearing change and no escalation
decision proceeds without appropriate independent review. The Sponsor-Investigator
(ChemicalQDevice, Kevin Kawchak, Chief Executive Officer) holds the combined IND and
IDE; the conflict of interest arising from the Sponsor-Investigator's relationship
to the developer of the Physical AI system is disclosed in full under 21 CFR part 54
and is mitigated structurally by this independent oversight, so that safety and
escalation decisions are not made by a financially interested party alone.
Table~\ref{tab:oversight-bodies} summarizes the oversight bodies, their cadence,
and their authority, and Figure~1 renders the governance structure.

\begin{table}[htbp]
\centering
\footnotesize
\caption{Three-tier independent safety oversight and the human-subjects-protection
structure for the combined IND and IDE study. Each body's review cadence and the
authority it exercises are stated; the binding halt rules drive the continue,
modify, pause, or halt decision.}
\label{tab:oversight-bodies}
\begin{tabularx}{\textwidth}{L{3.4cm} L{3.2cm} Y}
\toprule
\textbf{Body} & \textbf{Review cadence} & \textbf{Role and authority} \\
\midrule
Sponsor-Investigator (ChemicalQDevice, K. Kawchak) & Continuous & Holds combined IND and IDE; protocol custody, safety reporting under 21 CFR parts 312 and 812, and compliance assurance; conflict of interest disclosed and managed under 21 CFR part 54. \\
Reviewing IRB & Per submission and continuing review & Human-subjects protection; reviews protocol, consent, and the significant-risk device determination; may suspend or terminate approval. \\
Data Safety Monitoring Board (DSMB) & Each cohort boundary and any triggered review & Independent surgical, oncologic, biostatistical, and device-safety members; applies the prespecified halt rules; recommends continue, modify, pause, or halt; binding halt authority. \\
Independent Safety Monitor (ISM) & Per procedure, real-time & Case-level safety oversight independent of the operating team; authority to call an immediate stop. \\
Physical AI Safety Review Committee & Cadence of 90 days or less, plus ad hoc & Reviews telemetry and audit records, verification-before-generation gate-surface decisions, USL readiness, and software changes; recommends USL reassessment or oversight changes. \\
\bottomrule
\end{tabularx}
\end{table}

\figcaption{Table 1. The oversight bodies, their cadence, and their authority for
the combined IND and IDE first-in-human study.}

\begin{mermaidfig}
\node[mmgoal] (si)    at (4.2,0)      {Sponsor-Investigator\\ChemicalQDevice (K. Kawchak)\\COI disclosed, 21 CFR part 54};
\node[mmproc] (irb)   at (-0.6,-2.5)  {Reviewing IRB\\(human-subjects protection)};
\node[mmproc] (dsmb)  at (4.2,-2.5)   {DSMB\\cohort-boundary + triggered;\\halt authority};
\node[mmproc] (ism)   at (9.0,-2.5)   {Independent Safety Monitor\\per-procedure, real-time;\\stop authority};
\node[mmproc] (pais)  at (4.2,-5.2)   {Physical AI Safety Review\\Committee ($\leq$ 90-day cadence)};
\node[mmdark] (halt)  at (9.0,-5.2)   {Binding halt rules:\\device SAE $\geq$ 1 of 3 in a cohort\\OR 2 device-related deaths};
\node[mmdec]  (dec)   at (6.6,-7.6)   {Continue, modify,\\pause, or halt};
\draw[mmedge] (si) -- (irb);
\draw[mmedge] (si) -- (dsmb);
\draw[mmedge] (si) -- (ism);
\draw[mmedge] (si.south) to[out=-90,in=90,looseness=0.6] (pais.north);
\draw[mmedge] (dsmb) -- (halt);
\draw[mmedge] (halt) -- (dec);
\draw[mmedge] (ism.south) to[out=-90,in=60,looseness=0.7] (dec.north east);
\draw[mmedge] (pais.south) to[out=-90,in=180,looseness=0.7] (dec.west);
\end{mermaidfig}
\figcaption{Figure 1. Study governance and three-tier safety oversight. The
Sponsor-Investigator (ChemicalQDevice, Kevin Kawchak, with the conflict of interest
disclosed under 21 CFR part 54) operates under a reviewing IRB and three independent
bodies: the Data Safety Monitoring Board (cohort-boundary and triggered reviews,
with halt authority), the Independent Safety Monitor (per-procedure, real-time, with
stop authority), and the Physical AI Safety Review Committee (a cadence of 90 days or
less). The binding halt rules, a device-related serious adverse event in at least
one of three within a cohort or two device-related deaths at any point, drive the
continue, modify, pause, or halt decision.}

\subsubsection{Binding Halt Rules and Stopping Authority}

Two prespecified, binding halt rules govern the study and are applied by the DSMB.
First, a device-related or procedure-related serious adverse event in at least one
of three participants within a dose cohort triggers a mandatory DSMB review with
authority to halt. Second, two device-related deaths at any point in the study
trigger a mandatory DSMB review with authority to halt. On invocation of either
rule, enrollment and any patient-facing robotic activity stop pending review. These
rules operate independently of, and in addition to, the dose-finding rules of the
3+3 escalation (dose levels 160, 220, and 300 mg once daily, with a twenty-eight-day
dose-limiting-toxicity window) and the per-procedure stop authority of the
Independent Safety Monitor. The cyber-physical safety envelope provides further
hard stops that are not discretionary: force caps of no more than 3 N per arm and no
more than 18 N cumulative across the eight arms, a 10 kHz heartbeat with a
sixty-four-byte frame and a one-hundred-microsecond watchdog, arm-park within fifty
microseconds, a cross-arm emergency stop within three milliseconds to a commanded
0.0 N, and a hardware emergency stop within five hundred milliseconds consistent
with 21 CFR \S 312.404. The vascular no-fly gate samples at 10 kHz and enforces a
soft-warning advisory zone at 3.0 mm, blocking no-fly margins of 1.0 mm at the
superior mesenteric vein and portal vein and 1.5 mm at the hepatic artery, celiac
artery, and superior mesenteric artery, and a hard-stop emergency-stop zone at 5.0
mm. Any breach of a hard cyber-physical limit, a failure of the
verification-before-generation gate surface, or a loss of audit integrity
constitutes a clinical-hold ground for the autonomy-bearing system, on which all
patient-facing robotic activity ceases and the decommissioning and data-preservation
procedure is executed.

\subsubsection{Cybersecurity and System Integrity}

Cybersecurity is treated as a patient-safety control rather than an administrative
adjunct. The LLM-directed system is deployed on premises and isolated outside the
vendor kinematic stack, with no requirement for external network connectivity during
a procedure; the advisory model is hash-pinned to a deterministic seed and a
repository commit so that the validated configuration in use is exactly the
configuration that was tested and recorded. Access to identifiable records and to
the cyber-physical control surface is governed by a deny-by-default authorization
model, under which no role can read, write, or export a record except where an
access right has been explicitly granted, and electronic signatures are bound to the
records they authenticate. Every access, change, and export is written to a
hash-chained audit trail in which each entry cryptographically incorporates the hash
of the prior entry, so that the record is tamper-evident and any alteration is
detectable. A cybersecurity incident is one of the six Physical AI reporting
triggers under 21 CFR \S 312.32(g) and is reportable within fifteen calendar days,
or within seven calendar days where there is potential for patient harm; the audit
trail is preserved from twenty-four hours before to seventy-two hours after each
Physical AI adverse event under 21 CFR part 11. The other Physical AI triggers,
namely a serious Physical AI adverse event (seven or fifteen days), a system-drug
interaction (fifteen days), sustained AI model degradation across at least three
consecutive procedures or twenty-four cumulative hours, a sim-to-real divergence
exceeding twice the tolerance (a trajectory error above 4 mm or a force error above
1.0 N), and a digital-twin discrepancy, are reported on the same Subpart J pathway.

\subsubsection{The Seven Physical AI Record Types and 21 CFR Part 11}

Seven Physical AI record types are captured and retained for this trial class, in
addition to the conventional clinical-trial records. Table~\ref{tab:record-types}
enumerates each record type, its content, and its primary regulatory anchor. All
seven are bound to a deterministic seed and a repository commit by the hash-chained
audit trail (record type seven), so that any reported behavior can be replayed
deterministically and any requested record reconstructed. All electronic records,
signatures, and the cyber-physical telemetry stream are maintained under 21 CFR part
11, with validated systems, secure time-stamped audit trails that do not obscure
prior entries, role-based access controls, and electronic signatures bound to the
records they authenticate. Records are retained for the period required by 21 CFR
\S 312.57, namely at least two years after a marketing application is approved for
the indication or, if no application is filed or approved, two years after the
investigation is discontinued and the FDA is notified; the raw, losslessly
reconstructable command and sensor tier is preserved so that a requested record can
be reconstructed on FDA request under 21 CFR \S 312.57. Data released for analysis,
deposition, or publication are de-identified by the HIPAA Safe Harbor method, with
removal of all eighteen categories of protected health information identifiers.

\begin{table}[htbp]
\centering
\footnotesize
\caption{The seven Physical AI record types captured and retained for the
LLM-directed robotic Whipple, in addition to the conventional clinical-trial
records. All records are retained under 21 CFR part 11 for the period required by 21
CFR \S 312.57.}
\label{tab:record-types}
\begin{tabularx}{\textwidth}{L{0.7cm} L{3.6cm} Y}
\toprule
\textbf{No.} & \textbf{Record type} & \textbf{Content and regulatory anchor} \\
\midrule
1 & LLM advisory record & Each prompt, the supplied context, and the advisory output of the second-opinion oracle, hash-pinned to seed and commit; advisory only, never control authority (21 CFR \S 312.21(e)). \\
2 & Robot command record & Each command issued to the eight arms (56 degrees of freedom), with timestamps for deterministic replay. \\
3 & Sensor and telemetry record & The cyber-physical stream from 640 sensor channels (80 per arm; 100 kHz force, 10 kHz other), tiered from compact summaries to the raw losslessly reconstructable level. \\
4 & Safety-limit-event record & Emergency-stop activations and latencies, positional and force-envelope excursions ($\leq$ 3 N per arm, $\leq$ 18 N cumulative), and vascular no-fly gating events (soft 3.0 mm, no-fly 1.0 to 1.5 mm, hard-stop 5.0 mm). \\
5 & Verification-before-generation gate-surface record & Each accept, block, or escalate decision of the ten-gate suite, aligned with H.R. 9510 (Verification Before Generation in Physical AI Oncology Trials Act of 2026). \\
6 & Human-override and intervention record & Each manual override, takeover by the backup operator, and stop invoked by the Independent Safety Monitor. \\
7 & Hash-chained audit trail & Cryptographically chained binding of records 1 to 6 to a deterministic seed and repository commit; tamper-evident; preserved from $-$24 h to $+$72 h around each Physical AI adverse event (21 CFR part 11). \\
\bottomrule
\end{tabularx}
\end{table}

\figcaption{Table 2. The seven Physical AI record types, their content, and their
regulatory anchors. Record type seven binds the preceding six to a deterministic
seed and repository commit so that any reported behavior can be replayed
deterministically.}

\subsubsection{Quality Assurance, Re-Validation, and Decommissioning}

Quality assurance spans the clinical and the cyber-physical domains. Before every
procedure, the pre-procedure safety matrix is verified, the USL readiness rating is
confirmed at or above the deployment threshold, and the emergency-stop, positional,
and force envelopes are checked. The Phase 0 simulation-validation gate, which must
be satisfied before any patient-facing robotic activity, requires a USL of at least
7.0, at least one thousand simulations across at least two independent frameworks
(for example Isaac, MuJoCo, Gazebo, or PyBullet), and a sim-to-real agreement of
less than 2 mm in trajectory and less than 0.5 N in tip force. After any software
update, model change, or configuration change to the LLM-directed system, a Phase 0
re-validation is performed and the USL is reassessed before any further
patient-facing use; no update reaches a participant without re-validation. On any
clinical hold or closure, enrollment and patient-facing robotic activity stop
immediately, affected participants are transitioned to appropriate clinical care,
and a controlled decommissioning and data-preservation procedure is executed under
which the robotic system is locked out of patient contact, the deterministic
configuration and repository commit are recorded, and the complete hash-chained
command and telemetry record is sealed and retained for the regulatory retention
period. The computational evidence supporting this assurance approach includes
prior in silico and digital-twin work by the author, namely a sixty-second PDAC
Whipple with daraxonrasib simulation, an AI digital-twin PDAC simulation, a
quantitative systems pharmacology metastatic-PDAC simulation with verification,
validation, and uncertainty quantification, a one-hundred-thousand-patient in silico
five-arm Phase 3 PDAC study in triplicate, and end-to-end PDAC digital-twin trial
proposals \cite{qspmetpancre,fdadigtwinpc,chatgpt100kp,pdacdigtwinp}, together with a
sixteen-LLM code-generation evaluation against FDA adverse-event reporting data
\cite{llmcomp16fda}.

\subsection{Selected References}

The following selected references support the additional information in this section
and are cited in full in the back matter. The FDA Phase 1 guidance frames the
content and graded rigor appropriate to an original first-in-human IND
\cite{phase1guidance}, and the daraxonrasib and on-premises robotic Whipple sources
establish the drug mechanism and the device context \cite{daraxwhipple,onpremwhippl}.
The Physical AI regulatory adaptations of 21 CFR parts 312 and 50 and of ICH
guidance govern the autonomy-bearing overlay, the opt-out, and the Subpart J
reporting and record-retention controls \cite{cfr312paiuni,cfr050paiuni,ichpaistduni}.
The Verification Before Generation in Physical AI Oncology Trials Act of 2026 (H.R.
9510) anchors the verification-before-generation gate surface
\cite{congress9510}, and the trust, accountability, and site-readiness literature for
clinical Physical AI informs the governance and oversight design
\cite{clintrustpai,paiautospons,natlpaiplatf,paisitedocpk}. The PDAC epidemiology and
the prior computational evidence supporting the deferral and assurance rationale are
drawn from the indication and digital-twin sources
\cite{pdac060s2030,qspmetpancre,fdadigtwinpc,chatgpt100kp,pdacdigtwinp,llmcomp16fda,paipredict4x}.
A complete reference list appears in the back matter.
